\begin{table}[H]
\centering
\caption{Attrition --- Liberia}
\label{tab:lib_attrition}
\begin{threeparttable}
\begin{tabular}[t]{lcccc}
\toprule
\multicolumn{1}{c}{ } & \multicolumn{4}{c}{Experiment (Liberia)} \\
\cmidrule(l{3pt}r{3pt}){2-5}
\multicolumn{1}{c}{ } & \multicolumn{1}{c}{(1)} & \multicolumn{1}{c}{(2)} & \multicolumn{1}{c}{(3)} & \multicolumn{1}{c}{(4)} \\
\cmidrule(l{3pt}r{3pt}){2-2} \cmidrule(l{3pt}r{3pt}){3-3} \cmidrule(l{3pt}r{3pt}){4-4} \cmidrule(l{3pt}r{3pt}){5-5}
Sample & T missing & C missing & Adj.\ diff. & N \\
\midrule
Overall & 0.340 & 0.341 &   0.000 & 4784 \\
 & & & (  0.062) & \\
Grade 1 & 0.328 & 0.249 &   0.079 & 1184 \\
 & & & (  0.093) & \\
Grade 2 & 0.334 & 0.301 &   0.033 & 1233 \\
 & & & (  0.115) & \\
Grade 3 & 0.322 & 0.479 &  -0.157 & 1243 \\
 & & & (  0.108) & \\
Grade 4 & 0.379 & 0.316 &   0.063 & 1124 \\
 & & & (  0.109) & \\
\bottomrule
\end{tabular}
\begin{tablenotes}[para]
\item Notes: Attrition = missing endline. T missing and C missing are unadjusted arm means. Adjusted difference is the treatment coefficient from an OLS regression with strata FE; SEs clustered at the school grade group level.
\end{tablenotes}
\end{threeparttable}
\end{table}
